\chapter{Built-in benchmarks}
\label{ch:benchmarks}

\code{problem_name} selects the builder.
The physical numbers come from the dataclass instantiated at the bottom of \file{user_input.py}.
The solver does not mix parameters across benchmarks.
Changing \code{g}, a domain limit, or a boundary weight means editing that dataclass, not \code{USER_PINN_INPUTS}.

All six benchmarks define an exact or reference solution, a soft penalty, and a hard output map.
Table~\ref{tab:problems} is the index.
The second coordinate of the steady problems is \(y\).

\begin{table}[ht]
\centering
\caption{Built-in benchmarks. The field count is the number of network towers.}
\label{tab:problems}
\begin{tabular}{@{}llll@{}}
\toprule
\code{problem_name} & Fields & Second axis & Default idea \\
\midrule
\code{tidal_test} & 2 & time & linear tide, exact solution \\
\code{dam_break} & 2 & time & wet-bed Stoker problem \\
\code{transcritical_bump} & 2 & dummy time & steady jump over a bump \\
\code{steady_rotation} & 3 & \(y\) & solid-body rotation \\
\code{steady_vortex} & 3 & \(y\) & gradient-wind vortex \\
\code{steady_eddies} & 3 & \(y\) & Fourier eddies on a carved bed \\
\bottomrule
\end{tabular}
\end{table}

\section{Linear tidal test}

\code{problem_name = "tidal_test"} builds a linear tide in a flat basin of length \(L\), undisturbed depth \(H\), and linear friction \(r\).
The free-surface amplitude at the forced end is \(A\), and the period is \(T_{\mathrm{tide}}\) (the default is 12.42~hours).
The trainer works in nondimensional \(x\in[0,1]\) and \(\tau\) from 0 to \code{tau_max}, with \(\tau\) in units of the tidal period.
The default window is three tidal periods.
The comparison archives convert elevation back by \(A\), velocity by \(A L/(H T_{\mathrm{tide}})\), \(x\) to kilometers, and time to hours.

The residuals are the linear continuity and momentum balances,
\[
\eta_{\tau} + u_{x} = 0,
\qquad
u_{\tau} + \alpha \eta_{x} + \beta u = 0,
\]
with \(\alpha = g H T_{\mathrm{tide}}^{2}/L^{2}\) and \(\beta = r T_{\mathrm{tide}}\).
The exact solution is the complex wave number that satisfies those equations with the basin boundary data.

Hard mode matches that exact solution on \(\tau=0\), on \(x=0\), and on \(x=1\).
The interior is the transfinite blend of those three traces plus \(\tau\, x\,(1-x)\) times the network.
The collocation loss is the PDE residual.
The exact interior solution is not pasted into the domain.

Soft mode penalizes the network against the exact initial elevation and velocity, and against the exact elevation and velocity on both ends.
The default weights, used when the dataclass does not override them, are 20 on each initial field and 10 on each boundary field.
Sample counts default to 2048 initial points and 1024 points on each end, drawn uniformly at random on each penalty evaluation.

\section{Wet-bed dam break}

\code{problem_name = "dam_break"} is the scaled Stoker problem used by the reference dam-break study:
\(g=9.81\,\mathrm{m\,s^{-2}}\), left depth \(0.5\,\mathrm{m}\), right depth \(0.25\,\mathrm{m}\), still water, \(x\in[0,1]\,\mathrm{m}\), \(t\in[0,0.2]\,\mathrm{s}\), dam at \(x=0.5\,\mathrm{m}\).
The reference solution is the wet-bed Riemann solution.

The residuals are the one-dimensional conservation laws
\[
h_{t} + (hu)_{x} = 0,
\qquad
(hu)_{t} + \big(hu^{2} + \tfrac12 g h^{2}\big)_{x} = 0.
\]

\code{transition_width = 0} keeps the sharp initial step.
A positive width replaces the step by a \(\tanh\) jump of that scale, in meters.

Hard mode sets
\[
h = h_{0}(x) + \hat{t}\,\hat{x}\,(1-\hat{x})\, h_{\mathrm{net}},
\qquad
u = u_{0} + \hat{t}\,\hat{x}\,(1-\hat{x})\, u_{\mathrm{net}},
\]
where hats denote coordinates scaled to \([0,1]\) and \(h_{0}\) is the initial step.
At \(t=0\) and on both walls the network contribution vanishes, so the depth on each wall stays at the initial depth of that wall and the velocity stays at \(u_{0}\).

Soft mode penalizes the initial step and the two walls separately.
The default loss is
\[
L_{m} + L_{p} + 50\,\mathrm{MSE}(h_{\mathrm{IC}}) + 50\,\mathrm{MSE}(u_{\mathrm{IC}})
+ 10\,\mathrm{MSE}(h_{\mathrm{BC}}) + 10\,\mathrm{MSE}(u_{\mathrm{BC}}).
\]
The default 41 initial samples are packed toward the dam.
Each wall has 10 samples, uniform in time.
The network sees time in the first column and \(x\) in the second.
The other benchmarks pass \(x\) first.
You do not set that order; it is part of the dam-break builder.

\section{Transcritical bump}

\code{problem_name = "transcritical_bump"} is steady, frictionless flow over a parabolic bump with a stationary hydraulic jump.
The channel is \(x\in[0,25]\,\mathrm{m}\).
The bed is \(z = 0.2 - 0.05(x-10)^{2}\) for \(8<x<12\) and zero elsewhere.
The upstream unit discharge is \(q=0.18\,\mathrm{m^{2}\,s^{-1}}\) and the downstream depth is \(h=0.33\,\mathrm{m}\).
The reference solution follows from constant discharge and the specific-energy relation on each smooth branch.
The jump is placed where the Rankine--Hugoniot momentum condition changes sign.

The residuals are the steady balances with bed slope and no friction,
\[
(hu)_{x} = 0,
\qquad
\big(hu^{2} + \tfrac12 g h^{2}\big)_{x} + g h z_{x} = 0.
\]
The trainer still carries a time coordinate, on the dummy interval \([0,1]\).
The network is always evaluated at \(t=0\), so the learned field does not depend on time.

Hard mode enforces the mixed physical data exactly:
downstream depth is \(h(L)=h_{\mathrm{down}}\) for any network output, and upstream discharge is \(h(0)u(0)=q_{\mathrm{in}}\).
Upstream depth and downstream velocity remain free.
Soft mode penalizes those two conditions, each with default weight 50, and does not prescribe both \(h\) and \(u\) at each end.

\section{Steady solid-body rotation}

\code{problem_name = "steady_rotation"} is a steady square flow,
\[
u = -\omega y, \qquad v = \omega x,
\qquad
h = h_{\mathrm{center}} + \frac{\omega^{2}(x^{2}+y^{2})}{2g},
\]
on a flat frictionless bed.
The default square is \([-1,1]\times[-1,1]\), with \(\omega=1\,\mathrm{s^{-1}}\) and \(h_{\mathrm{center}}=1\,\mathrm{m}\).
The residuals are the steady conservative shallow-water equations with no friction and no bed slope.
There is no time axis.
\code{trial_nx} and \code{trial_nt} are the \(x\) and \(y\) station counts.

Hard mode returns the exact field plus a bubble \(s_{x}(1-s_{x})s_{y}(1-s_{y})\) times the network, where \(s_{x}\) and \(s_{y}\) are the coordinates scaled to \([0,1]\).
The bubble is zero on the boundary, so the boundary values are exact.
A zero network is already the exact solution everywhere, because the exact field is part of the output.
This benchmark is a useful wiring check.
It is a weak test of whether a family can represent an unknown interior.

Soft mode penalizes \(h\), \(u\), and \(v\) on the four edges against the exact traces.
The default edge has 32 stations, and each field weight defaults to 50.

\section{Steady vortex}

\code{problem_name = "steady_vortex"} is an isolated vortex in gradient-wind balance on a flat bed.
The azimuthal speed is \(\omega r \exp(-r^{2}/(2\sigma^{2}))\), the radial velocity is zero, and the free surface satisfies
\[
\frac{v_{\theta}^{2}}{r} + f v_{\theta} = g \frac{\partial h}{\partial r}.
\]
The default Coriolis parameter is \(f=2\,\mathrm{s^{-1}}\), a rotating-table value rather than the planetary one.
Set \(f=0\) on \code{SteadyVortexParameters} for a non-rotating bucket vortex.
The default core rate is \(\omega=4\,\mathrm{s^{-1}}\), the speed peaks at \(r=\sigma=0.45\,\mathrm{m}\), and the central depth is \(0.40\,\mathrm{m}\).
The square is \([-1.5,1.5]\times[-1.5,1.5]\), several vortex radii across, so the boundary velocity is a small fraction of the peak.
Continuity is satisfied identically.
There is no manufactured body force.

Hard mode lifts the boundary traces with a Coons patch and adds the same boundary-vanishing bubble times the network.
The exact interior solution is not inserted.
Soft mode penalizes all three fields on the four edges against the exact boundary values, with the same default station count and weights as the rotation problem.

\section{Steady eddy lattice}

\code{problem_name = "steady_eddies"} is the benchmark selected by the shipped \file{user_input.py}.
The stream function is a degenerate Fourier pair
\[
\psi = A \big[
\sin(m\pi x/L)\sin(n\pi y/L)
+ \sin(n\pi x/L)\sin(m\pi y/L)
\big],
\]
with velocity \((\psi_{y},-\psi_{x})\) and depth \(H_{0}+\gamma\psi\).
Both mode pairs share the eigenvalue \(\mu=(m^{2}+n^{2})\pi^{2}/L^{2}\).
The free surface that balances advection is
\[
\eta = C - \frac{|\mathbf{u}|^{2}}{2g} - \frac{\mu\psi^{2}}{2g},
\]
and the bed is \(z=\eta-h\).
Continuity and momentum hold identically.
There is no manufactured body force and no shock.
The walls are slip walls: depth equals \(H_{0}\) and the normal velocity vanishes.
The defaults are \(A=0.015\,\mathrm{m^{2}\,s^{-1}}\), \(H_{0}=0.80\,\mathrm{m}\), \(\gamma=4\), modes \((m,n)=(7,4)\), Bernoulli constant \(C=1\), and the unit square.
Keep \(C\) large enough that the bed stays below the free surface.
Mode~7 needs more than twenty stations across the square; the shipped trial lattice uses 48.

Hard mode is the Coons lift of the boundary traces plus the bubble times the network.
The exact interior is not inserted.
Soft mode enforces the slip condition only: depth equals \(h_{\mathrm{mean}}\) on the walls, and the normal velocity vanishes.
Tangential speed is left free so that penalty does not dominate the Optuna score.
\code{soft_w_u} and \code{soft_w_v} are averaged into the normal-velocity weight.

\section{Choosing a constraint mode}

Use hard mode when the boundary traces should be exact for every set of weights, and the loss should be the PDE residual alone.
That is the shipped choice for the tide and for the eddy lattice.

Use soft mode when the boundary data are mixed, or when you want the penalty weights to be part of the experiment.
The dam-break reference and the bump preset use soft mode.
The bump's hard map already enforces the mixed discharge and depth conditions, so hard mode is also a valid bump run.
Give the bump a new \code{study_name} if you change the constraint mode of an existing study; the trials are not comparable across modes.

Use a new database whenever you change \code{problem_name}.
Trials from the tide and trials from the dam break must not share a study.
The loss values are residuals of different equations.

\section{Adding a benchmark}
\label{sec:new-problem}

A new case is a builder that returns a \code{Problem} from \file{swe_pinn/physics/common.py}: a name, a \code{Domain1D}, a residual function, an output function, and, when you want soft mode or an analytic comparison, a soft-penalty function and an exact solution.
Register the name in \code{build_problem} in \file{swe_pinn/problems/__init__.py}, and add a dataclass to \file{user_input.py}.
The residual returns two tensors for a one-dimensional problem or three tensors for a steady \((x,y)\) problem.
Set \code{n_outputs} to 2 or 3 to match.
The trainer, the three families, and the Optuna loop stay as they are.
Point \code{storage_url} and \code{study_name} at a new study before the first trial.
